%==============================================================================
\section{How far can the two-phase heuristic be from optimal?}\label{sec:gap}
%==============================================================================

The two-phase heuristic of Section~\ref{ssec:jgp} is the standard practical method
for the SSP. Where it has been measured against exact optima its suboptimality is
small \citep{Burger2015}. No worst-case analysis of it has been published. This
section gives one.

The source of any gap is easy to state. The heuristic commits to a grouping with the
\emph{fewest} groups before it sequences anything, whereas the optimal schedule may
be served by a larger grouping with cheaper transitions. Everything below makes that
statement precise, bounds the resulting loss, and identifies when it is zero.

\paragraph{Which value is being bounded.}
Section~\ref{ssec:frames} distinguished two values: the walk value $\Hwalk$, and the
value $H$ that the method returns. Every result below names the one it is about. The
distinction is not a technicality. Section~\ref{ssec:ktnsgap} shows that the
canonical example family of this problem has a positive walk gap and a zero
heuristic gap, so a result proved about $\Hwalk$ and reported as a result about the
heuristic would be false.

\paragraph{How the results were checked.}
Every statement in this section was verified computationally by a program written
independently of the solvers used elsewhere in this report, described in
Appendix~\ref{app:verify}. Statements labelled Theorem, Proposition, Lemma or
Corollary also have proofs, given here or in Appendix~\ref{app:proofs}. Statements
labelled Observation were established by computation only. Statements labelled
Conjecture have neither a proof nor a claim to one.

\subsection{The gap is the price of using the fewest groups}\label{ssec:exact}

A grouping is given a cost by ordering its configurations as cheaply as possible.

\begin{definition}[Feasible grouping and its cost]\label{def:grouping}
A \emph{feasible grouping} is a partition $\mathcal{P}=\rset{G_1,\dots,G_p}$ of $J$
into feasible groups, together with a configuration $C_l\supseteq\bigcup_{j\in G_l}T_j$
for each group. Its cost is the cheapest free-initial walk through its
configurations,
\[
  \gamma(\mathcal{P})\;=\;\min_{\sigma\in S_p}\ \sum_{l=1}^{p-1}
  d\bigl(C_{\sigma(l)},C_{\sigma(l+1)}\bigr).
\]
\end{definition}

The Job Grouping Problem seeks a feasible grouping of minimum cardinality $\Kst$,
and $\Hwalk$ is $\gamma$ minimised over those. Removing the cardinality restriction
makes the framework exact.

\begin{theorem}[Grouping exactness]\label{thm:grouping}
$\displaystyle \Zst=\min_{\mathcal{P}}\gamma(\mathcal{P})$, where the minimum runs
over \emph{all} feasible groupings, of every cardinality.
\end{theorem}

The proof is in Appendix~\ref{app:proofs}. It is a pair of constructions. Any
feasible grouping, processed group by group, is a schedule of cost
$\gamma(\mathcal{P})$, which gives one inequality. Any optimal schedule, with its
jobs grouped by the magazine under which they run, is a feasible grouping of no
greater cost, which gives the other.

\begin{corollary}[The walk gap]\label{cor:gap}
\[
  \Hwalk-\Zst \;=\; \min_{\lvert\mathcal{P}\rvert=\Kst}\gamma(\mathcal{P})
  \;-\;\min_{\mathcal{P}}\gamma(\mathcal{P})\;\ge\;0,
\]
and it is strictly positive exactly when no minimum-cardinality grouping attains the
overall minimum of $\gamma$.
\end{corollary}

This locates the loss exactly. It is not a failure of the sequencing phase, which
orders whatever it is given optimally. It is the price of the cardinality
restriction imposed by the grouping phase.

\begin{example}[Where the ring loses, in the walk model]\label{ex:ring-why}
The optimum of the $6$-ring is realised by a \emph{four}-configuration grouping,
\[
  \rset{1,2,3}\ \to\ \rset{1,3,4}\ \to\ \rset{1,4,5}\ \to\ \rset{1,5,6},
\]
in which consecutive configurations differ in a single tool, for $\gamma=3$. Every
minimum-cardinality grouping uses three configurations and costs $4$
(Example~\ref{ex:frames-ring}). The walk gap of one is exactly the difference in
Corollary~\ref{cor:gap}: the cheaper grouping is not of minimum cardinality, so the
grouping phase never offers it to the sequencer.
\end{example}

\begin{figure}[t]\centering
\begin{tikzpicture}[font=\footnotesize,>=Stealth,
   cfg/.style={draw=cPurple!70,semithick,fill=white,rounded corners,minimum width=1.6cm,minimum height=0.85cm},
   bad/.style={draw=cRed!70,semithick,fill=white,rounded corners,minimum width=1.75cm,minimum height=0.85cm},
   node distance=1.15cm]
 \begin{scope}
  \node[font=\footnotesize\bfseries,anchor=west] at (-0.9,1.15)
     {four configurations: cost $3$};
  \node[cfg] (c1) {$\{1,2,3\}$};
  \node[cfg,right=of c1] (c2) {$\{1,3,4\}$};
  \node[cfg,right=of c2] (c3) {$\{1,4,5\}$};
  \node[cfg,right=of c3] (c4) {$\{1,5,6\}$};
  \draw[->,semithick,black!60] (c1)--(c2) node[midway,above,font=\scriptsize]{$+4$};
  \draw[->,semithick,black!60] (c2)--(c3) node[midway,above,font=\scriptsize]{$+5$};
  \draw[->,semithick,black!60] (c3)--(c4) node[midway,above,font=\scriptsize]{$+6$};
  \node[below=0.5mm of c1,font=\scriptsize,text=black!60] {serves $1,2$};
  \node[below=0.5mm of c2,font=\scriptsize,text=black!60] {serves $3$};
  \node[below=0.5mm of c3,font=\scriptsize,text=black!60] {serves $4$};
  \node[below=0.5mm of c4,font=\scriptsize,text=black!60] {serves $5,6$};
 \end{scope}
 \begin{scope}[yshift=-3.2cm]
  \node[font=\footnotesize\bfseries,anchor=west] at (-0.9,1.15)
     {three configurations: cost $4$};
  \node[bad] (a) at (0.6,0) {$\{1,2,3\}$};
  \node[bad] (b) at (3.9,0) {$\{3,4,5\}$};
  \node[bad] (c) at (7.2,0) {$\{5,6,1\}$};
  \draw[->,semithick,black!60] (a)--(b) node[midway,above,font=\scriptsize]{$+4,+5$};
  \draw[->,semithick,black!60] (b)--(c) node[midway,above,font=\scriptsize]{$+6,+1$};
  \node[below=0.5mm of a,font=\scriptsize,text=black!60] {serves $1,2$};
  \node[below=0.5mm of b,font=\scriptsize,text=black!60] {serves $3,4$};
  \node[below=0.5mm of c,font=\scriptsize,text=black!60] {serves $5,6$};
 \end{scope}
\end{tikzpicture}
\caption{The $6$-ring in configuration space, in the walk model. Above: the optimal
schedule visits four configurations, each one tool from the next, for a cost of
three. Below: the heuristic may use only the $\Kst=3$ minimum-cardinality groups, and
every pair of the resulting configurations differs in two tools, so the best it can
do in this model is four. The extra unit is one re-insertion of tool~1, which the
middle configuration cannot hold. Figure~\ref{fig:frames} shows what happens when
the KTNS step is applied instead.}
\label{fig:gtsp}
\end{figure}

\subsection{Zero-gap classes and worst-case guarantees}\label{ssec:ratio}

This subsection proves bounds that hold for \emph{every} instance and every capacity.
Two quantities recur. The first is the tool excess $q=\lvert\U\rvert-b$ of
Section~\ref{ssec:lb}. The second is the number of \emph{re-insertions} a schedule
makes: insertions of tools that were in the magazine earlier and were removed. Write
$R_{\mathrm{opt}}$ for the re-insertions made by an optimal schedule. Throughout,
$\Kst\ge2$, since at $\Kst=1$ the heuristic is trivially exact, and all costs are
free-initial.

\begin{lemma}[Cost identity]\label{lem:costid}
Restrict every magazine state to a subset of $\U$, which is possible whenever
$\Kst\ge2$ and never increases cost. The free-initial cost of any full-magazine
schedule then equals $q+R$, where $R\ge0$ counts re-insertions.
\end{lemma}

\begin{proof}
Each of the $q$ tools of $\U$ outside the first fill is required by some job, so it
is inserted at least once. Split all insertions into first-time insertions, of which
there are exactly $q$, and the rest, which are by definition the $R$ re-insertions.
\end{proof}

Hence $\Zst=q+R_{\mathrm{opt}}$, and writing $R_{\mathrm{walk}}$ for the walk model's
re-insertions, $\Hwalk-\Zst=R_{\mathrm{walk}}-R_{\mathrm{opt}}$. The walk gap is
exactly the excess re-insertions forced by the cardinality restriction. On the
$6$-ring the single extra insertion is one such re-insertion: tool~1, needed by the
first and last groups but not the middle one, is removed at a group boundary and
reloaded (Figure~\ref{fig:frames}).

\begin{lemma}[Transition cap]\label{lem:transcap}
In any ordering of configurations contained in $\U$, every transition inserts at most
$\min(b,q)$ tools. Hence $\Hwalk\le(\Kst-1)\min(b,q)$.
\end{lemma}

\begin{proof}
A transition into $C_{k+1}$ inserts $\lvert C_{k+1}\setminus C_k\rvert$ tools. This
is at most $\lvert C_{k+1}\rvert=b$, and also at most
$\lvert\U\setminus C_k\rvert=q$, since $C_{k+1}\subseteq\U$ and
$\lvert C_k\rvert=b$. A covering configuration inside $\U$ exists for every group, so
some minimum-cardinality grouping uses only such configurations, and $\Hwalk$ is at
most its cost.
\end{proof}

These two lemmas give the section's central guarantee. Lemma~\ref{lem:transcap} caps
the heuristic's absolute cost, Corollary~\ref{cor:lb} floors the optimum, and
dividing one by the other bounds how far from optimal the method can stray.

\begin{theorem}[Worst-case guarantee]\label{thm:uncond}
For every instance with $\Kst\ge2$,
\[
  \frac{H}{\Zst}\ \le\ \frac{\Hwalk}{\Zst}
  \ \le\ \frac{(\Kst-1)\min(b,q)}{\max(\Kst-1,\ q)}
  \ \le\ \min\bigl(b,\ \Kst-1,\ q\bigr).
\]
\end{theorem}

\begin{proof}
The first inequality is Proposition~\ref{prop:frames}. For the second, divide
Lemma~\ref{lem:transcap} by $\Zst\ge\max(\Kst-1,q)\ge1$ (Corollary~\ref{cor:lb}). For
the third, $(\Kst-1)q/\max(\Kst-1,q)=\min(\Kst-1,q)$ together with $\min(b,q)\le q$
gives the $\min(\Kst-1,q)$ term, while $(\Kst-1)/\max(\Kst-1,q)\le1$ together with
$\min(b,q)\le b$ gives the $b$ term.
\end{proof}

This subsumes the trivial $b$-approximation that follows from
$\Hwalk\le b(\Kst-1)$ and $\Zst\ge\Kst-1$, which was the bound the internship began
with. It tightens automatically on instances with few groups, with a small tool
excess, or with a small magazine.

\begin{corollary}[Zero-gap classes, every $b$]\label{cor:zerogap}
\emph{(i)} If $\Kst=2$ the gap is zero, and indeed $H=\Hwalk=\Zst=q$.
\emph{(ii)} If $\lvert\U\rvert\le b+1$ the gap is zero.
\emph{(iii)} If $\Kst=3$ then $H/\Zst\le2$.
\end{corollary}

\begin{proof}
All three are instances of Theorem~\ref{thm:uncond}, whose ratio bound is
$\Kst-1=1$, $q\le1$ and $\Kst-1=2$ respectively. For the equality in (i), pad the
first group's configuration to size $b$ using tools from the second group's
requirement, and pad the second from the first. The single transition then inserts
exactly the $q$ tools of $\U$ missing from the first configuration, so
$\Hwalk\le q\le\Zst\le H\le\Hwalk$ by the coverage bound.
\end{proof}

A positive gap therefore requires $\Kst\ge3$.

\begin{corollary}[Small optima are gap-free]\label{cor:smallZ}
If $\Zst\le2$ then the gap is zero, for every $b$ and every $\Kst$.
\end{corollary}

\begin{proof}
$\Zst\ge q$, so $q\le2$. If $q\le1$ the gap vanishes by
Corollary~\ref{cor:zerogap}(ii). If $q=2$ then $\Zst=q$, and the walk argument of
Appendix~\ref{app:proofs} yields a walk of at most $q+1=3$ configurations, each
serving a job, exhibiting a feasible grouping of that cardinality; for $\Kst\le2$ the
gap vanishes by Corollary~\ref{cor:zerogap}(i), and for $\Kst=3$ the walk grouping is
itself of minimum cardinality and has cost $\Zst$, forcing $\Hwalk=\Zst$.
\end{proof}

\begin{corollary}[$\Zst\le3$ forces a walk gap of at most one]\label{cor:z3}
If $\Zst\le3$ then $\Hwalk-\Zst\le1$, for every $b$ and every $\Kst$.
\end{corollary}

\begin{proof}
Write $\Zst=q+R_{\mathrm{opt}}$, so $q\le3$. If $q\le1$ or $\Kst\le2$ the gap is zero
by Corollary~\ref{cor:zerogap}. Assume $q\in\rset{2,3}$ and $\Kst\ge3$. The excised
optimal walk has $p\le\Zst+1\le4$ configurations, each serving a job, and exhibits a
feasible $p$-grouping, so $\Kst\le p\le4$. If $\Kst=4$ then $p=\Kst$, the walk's
grouping is of minimum cardinality and has cost $\Zst$, forcing $\Hwalk=\Zst$. There
remains $\Kst=3$. For $(q,R_{\mathrm{opt}})=(3,0)$ the merge argument in
Appendix~\ref{app:proofs} gives a gap of at most one. For
$(q,R_{\mathrm{opt}})=(2,1)$ the bound of Proposition~\ref{prop:k3} reads
$\max\bigl(0,\min(2,\lfloor(2b-2)/3\rfloor)-1\bigr)\le1$.
\end{proof}

\subsection{The extremal behaviour at three groups}\label{ssec:law}

Corollary~\ref{cor:zerogap} places every positive-gap instance at $\Kst\ge3$. This
subsection works out the case $\Kst=3$, which is where the extremes occur. All
results here are about the walk value; Section~\ref{ssec:ktnsgap} then asks what the
KTNS step does to them.

\begin{proposition}[Exact walk cost at $\Kst=3$]\label{prop:hk3}
If $\Kst=3$ then, with configurations restricted to subsets of $\U$,
\[
  \Hwalk\;=\;q\;+\;\min\ \min\bigl(x_{ab},\,x_{bc},\,x_{ca}\bigr),
  \qquad x_{uv}=\bigl\lvert(C_u\cap C_v)\setminus C_w\bigr\rvert,
\]
where the outer minimum runs over feasible $3$-groupings and choices of their
configurations, and $w$ denotes the third index in each case.
\end{proposition}

\begin{proof}
The three configurations of a feasible $3$-grouping are distinct, since two equal
ones would merge their groups and contradict $\Kst=3$, and together they cover $\U$.
For a path $(C_a,C_b,C_c)$, splitting the cost into first-time insertions and
re-insertions (Lemma~\ref{lem:costid}) gives cost
$q+\lvert(C_a\cap C_c)\setminus C_b\rvert$: the re-inserted tools are exactly those
of $C_c$ present in $C_a$ but absent from $C_b$. Minimising over the three possible
middle configurations selects the smallest of the three overlap terms, and $\Hwalk$
then minimises over groupings and configuration choices.
\end{proof}

This turns grouping selection at $\Kst=3$ into an explicit criterion: choose the
grouping whose configurations have the smallest pairwise overlap outside the third.
Section~\ref{ssec:grouping-sel} returns to it.

\begin{lemma}[Overlap counting]\label{lem:overlap}
Any three configurations $C_a,C_b,C_c\subseteq\U$ of size $b$ with
$C_a\cup C_b\cup C_c=\U$ satisfy
$\min(x_{ab},x_{bc},x_{ca})\le\lfloor(2b-q)/3\rfloor$.
\end{lemma}

\begin{proof}
Count tools by how many of the three configurations contain them. Let $a$ be the
number in exactly one, $\sum x:=x_{ab}+x_{bc}+x_{ca}$ the number in exactly two, and
$c_3$ the number in all three. Counting tools gives $a+\sum x+c_3=b+q$, and counting
memberships gives $a+2\sum x+3c_3=3b$. Subtracting, $\sum x=2b-q-2c_3\le2b-q$. The
minimum of three numbers is at most their mean, and it is an integer.
\end{proof}

\begin{proposition}[The walk gap at $\Kst=3$]\label{prop:k3}
Let $\Kst=3$ and $R_{\mathrm{opt}}=\Zst-q$. Then
\[
  \Hwalk-\Zst\ \le\ \max\Bigl(0,\ \min\bigl(q,\ \lfloor(2b-q)/3\rfloor\bigr)
  -R_{\mathrm{opt}}\Bigr).
\]
Consequently $\Hwalk-\Zst\le1$ holds
\emph{(i)} for every instance with $b\le4$;
\emph{(ii)} whenever $2b\le q+3R_{\mathrm{opt}}+5$;
\emph{(iii)} whenever $R_{\mathrm{opt}}\ge q-1$; and
\emph{(iv)} at $\Zst=q=3$ when two of the four walk groups have combined requirement
fitting one magazine.
The bound can therefore fail only in the corner
$2b\ge q+3R_{\mathrm{opt}}+6$ with $R_{\mathrm{opt}}\le q-2$, which forces $b\ge5$.
\end{proposition}

\begin{proof}
The configurations of any feasible $3$-grouping cover $\U$, so
Proposition~\ref{prop:hk3} with Lemma~\ref{lem:overlap} gives
$\Hwalk\le q+\lfloor(2b-q)/3\rfloor$, while Lemma~\ref{lem:transcap} gives
$\Hwalk\le2\min(b,q)\le 2q$. Subtracting $\Zst=q+R_{\mathrm{opt}}$ yields the bound.
For (i) at $b=3$, $\lfloor(6-q)/3\rfloor\le1$ for every $q\ge1$. At $b=4$ the bound
is at most $1$ for $q\ge3$; for $q=2$, either $R_{\mathrm{opt}}\ge1$ and the bound is
at most $1$, or $\Zst=q\le2$ and the gap vanishes by Corollary~\ref{cor:smallZ}; and
$q\le1$ is Corollary~\ref{cor:zerogap}(ii). Parts (ii) and (iii) are read off the
bound, and part (iv) is proved in Appendix~\ref{app:proofs}.
\end{proof}

\begin{proposition}[A bound for every $\Kst$]\label{prop:genk}
For every instance with $\Kst\ge2$,
\[
  \Hwalk-\Zst\ \le\ \max\Bigl(0,\
  \Bigl\lfloor\frac{q\bigl(b(\Kst-1)-q\bigr)}{b+q}\Bigr\rfloor-R_{\mathrm{opt}}\Bigr).
\]
\end{proposition}

\begin{proof}
Fix a minimum-cardinality grouping. Its configurations $C_1,\dots,C_{\Kst}\subseteq\U$
are pairwise distinct and cover $\U$. In an ordering $\sigma$ the path cost is
$q+\sum_t(\mathrm{runs}_\sigma(t)-1)$, where $\mathrm{runs}_\sigma(t)$ counts the
maximal blocks of consecutive configurations containing $t$, since every block after
the first begins with a re-insertion (Lemma~\ref{lem:costid}). Let $d_t$ be the number
of configurations containing $t$, so that $\sum_t d_t=\Kst b$ over $b+q$ tools. Under
a uniformly random ordering,
$\mathbb{E}[\mathrm{runs}(t)-1]=(d_t-1)(\Kst-d_t)/\Kst$, which is concave in $d_t$, so
by Jensen's inequality the expected total excess is at most
$q\bigl(b(\Kst-1)-q\bigr)/(b+q)$. Some ordering achieves at most the mean, the total is
an integer, and subtracting $\Zst=q+R_{\mathrm{opt}}$ gives the claim.
\end{proof}

At $\Kst=3$ this and Proposition~\ref{prop:k3} are incomparable and both are attained
at the $6$-ring. Proposition~\ref{prop:genk} grows linearly in $\Kst$ and is
unconditional.

The next result shows that the natural guess $\mathrm{gap}\le\Kst-2$, which held on
all the evidence collected before this internship, is false for the walk value. All
the earlier evidence lay at $b\le4$, which is exactly the region
Proposition~\ref{prop:k3} proves.

\begin{proposition}[The $\Kst-2$ bound fails for the walk value]\label{prop:refute}
There are instances with $\Kst=3$ and $\Hwalk-\Zst=2$. One, with $b=5$ and
$\U=\rset{1,\dots,9}$, is
\[
  W:\quad T=\bigl(\rset{1,2,3,5,8},\ \rset{1,5,6},\ \rset{2,6},\ \rset{3,7,9},\
  \rset{4,5,6,9}\bigr),
\]
with $q=4$, $\Zst=4$ and $\Hwalk=6$, both values obtained by exhaustive computation.
The instance lies exactly on the boundary of the regimes of
Proposition~\ref{prop:k3}, at $2b=q+3R_{\mathrm{opt}}+6$, and attains that
proposition's bound with equality.
\end{proposition}

\begin{observation}\label{obs:refute-ktns}
On the same instance $W$, the KTNS step gives $H=5$, so the heuristic gap is one and
the guess $\mathrm{gap}\le\Kst-2$ is not contradicted by it.
\end{observation}

\begin{conjecture}[The $4/3$ ratio at $b=3$]\label{conj:43}
For every instance with $b=3$, $H/\Zst\le4/3$.
\end{conjecture}

The evidence is a complete enumeration of two families. The first is every instance
in which each job requires exactly two tools, with $b=3$, at most six tools and at
most six jobs: $10{,}691$ instances, of which the maximum ratio is exactly $4/3$,
attained at the $6$-ring in the walk model. The second is every instance with $b=3$
whose jobs require two or three tools, with at most five tools and at most five jobs:
a further $6{,}065$ instances, none exceeding $4/3$.

\begin{proposition}[The $4/3$ bound holds whenever $\Kst\le3$]\label{prop:43route}
Every instance with $b=3$ and $\Kst\le3$ satisfies $\Hwalk/\Zst\le4/3$, and hence
$H/\Zst\le4/3$. Conjecture~\ref{conj:43} is therefore open only for $\Kst\ge4$.
\end{proposition}

\begin{proof}
For $\Kst\le2$ the gap is zero (Corollary~\ref{cor:zerogap}). For $\Kst=3$,
Proposition~\ref{prop:k3}(i) gives a walk gap of at most one, and a positive gap
forces $\Zst\ge3$ (Corollary~\ref{cor:smallZ}), so
$\Hwalk/\Zst=1+(\Hwalk-\Zst)/\Zst\le1+1/3$. Proposition~\ref{prop:frames} carries the
bound to $H$.
\end{proof}

The reduction to $\Kst\ge4$ is not vacuous. Positive-walk-gap instances with $b=3$
and $\Kst=4$ exist; one is
$T=(\rset{1,3,6},\rset{1,4,7},\rset{2,3},\rset{2,5},\rset{2,7})$, with $\Zst=4$ and
$\Hwalk=5$. No ratio above $4/3$ has been found in that regime.

\subsection{What the KTNS step does to the walk gap}\label{ssec:ktnsgap}

Everything in Section~\ref{ssec:law} bounds $\Hwalk$. By
Proposition~\ref{prop:frames} those bounds also hold for $H$. The question this
subsection asks is how much slack that inequality hides. The answer is: a great deal,
and on exactly the instances the literature uses as examples.

\begin{observation}[The ring family has no heuristic gap]\label{obs:rings}
For the $k$-ring at $b=3$ with $k=3,\dots,9$, the heuristic is optimal:
$H=\Zst$ in every case. The walk value is one above the optimum for
$k=6,\dots,9$. Table~\ref{tab:rings} gives the values.
\end{observation}

\begin{table}[h]\centering
\begin{tabular}{ccccccc}
\toprule
$k$ & $\Kst$ & $\Zst$ & $\Hwalk$ & walk gap & $H$ & heuristic gap\\
\midrule
3 & 1 & 0 & 0 & 0 & 0 & 0\\
4 & 2 & 1 & 1 & 0 & 1 & 0\\
5 & 3 & 2 & 2 & 0 & 2 & 0\\
6 & 3 & 3 & 4 & 1 & 3 & 0\\
7 & 4 & 4 & 5 & 1 & 4 & 0\\
8 & 4 & 5 & 6 & 1 & 5 & 0\\
9 & 5 & 6 & 7 & 1 & 6 & 0\\
\bottomrule
\end{tabular}
\caption{The $k$-ring at $b=3$, free-initial costs, computed exhaustively. The walk
value separates from the optimum at $k=6$ and stays one above it. The value the
heuristic actually returns equals the optimum for every $k$ tested. The mechanism is
the one drawn in Figure~\ref{fig:frames}.}\label{tab:rings}
\end{table}

This is why the distinction of Section~\ref{ssec:frames} cannot be dropped. The
$6$-ring is the standard example of a positive gap for this heuristic, and for the
heuristic as defined it has no gap at all. The same holds for the $\Kst=4$ instance
at the end of Section~\ref{ssec:law}, whose walk gap is one and whose heuristic gap
is zero, and for the witness $W$ of Proposition~\ref{prop:refute}, whose walk gap is
two and whose heuristic gap is one.

Positive heuristic gaps do exist. They are simply rarer and smaller than the walk
model suggests.

\begin{observation}[The smallest sub-optimal instance]\label{obs:smallest}
The instance with $b=4$, four jobs and $\U=\rset{1,\dots,7}$ given by
\[
  T=\bigl(\rset{1,2},\ \rset{3,4},\ \rset{1,3,5,6},\ \rset{2,5,6,7}\bigr)
\]
has $\Kst=3$, $\Zst=3$ and $H=4$, so the heuristic is suboptimal on it. An exhaustive
search over all instances with at most four jobs finds no smaller one.
\end{observation}

\begin{observation}[Positive heuristic gaps are rare and small]\label{obs:census}
In a random sample of $420$ instances, drawn with three to five jobs, four to eight
tools, capacity between two and five, and each job's requirement a uniformly random
subset of admissible size, exactly one has $H>\Zst$. Every positive heuristic gap
found anywhere in this study equals one, and every instance carrying one has
$b\ge4$.
\end{observation}

The last part of that observation is worth pausing on. A positive heuristic gap
requires the magazine to be large relative to the jobs. That is precisely the regime
in which \citet{Burger2015} report the decomposition degrading: as the difference
between the magazine size and the largest job size grows, the quality of the
heuristic's solutions decreases. Their finding is empirical and industrial, this one
is combinatorial and exhaustive, and they agree.

Unboundedness survives, but not with the family the walk model suggests.

\begin{proposition}[Unbounded walk gap]\label{prop:unbounded}
Let $R_g$ consist of $g$ tool-disjoint copies of the $6$-ring, so that $b=3$ and
$\lvert\U\rvert=6g$. Then $\Zst=6g-3$ and $\Hwalk=7g-3$, so the walk gap $g$ grows
without bound, while the walk ratio $(7g-3)/(6g-3)$ decreases from $4/3$ towards
$7/6$.
\end{proposition}

The proof is in Appendix~\ref{app:proofs}. It pairs the coverage bound with an
explicit schedule attaining it, and counts the heuristic's cost as $g$ within-copy
contributions of four and $g-1$ inter-copy transitions of three.

By Observation~\ref{obs:rings} this construction says nothing about $H$: the ring has
no heuristic gap, so copies of it have none either. Replacing the seed repairs the
construction.

\begin{observation}[Unbounded heuristic gap, with a different seed]\label{obs:unbounded-true}
Let $S$ be the four-job instance of Observation~\ref{obs:smallest}, which has $b=4$,
$\lvert\U\rvert=7$ and heuristic gap one, and let $S_g$ consist of $g$ tool-disjoint
copies of it. Exhaustive computation gives
\[
  \begin{array}{c|ccccc}
   g & \lvert\U\rvert & \Kst & \Zst & H & H-\Zst\\\hline
   1 & 7  & 3 & 3  & 4  & 1\\
   2 & 14 & 6 & 10 & 12 & 2
  \end{array}
\]
so $\Zst=7g-4=\lvert\U\rvert-b$ and the heuristic gap equals $g$ on both.
\end{observation}

\begin{conjecture}\label{conj:unbounded}
For every $g\ge1$ the instance $S_g$ of Observation~\ref{obs:unbounded-true} has
$\Zst=7g-4$ and $H-\Zst=g$. In particular the heuristic gap is unbounded.
\end{conjecture}

\begin{support}
The optimum is immediate: the coverage bound gives $\Zst\ge7g-4$, and processing the
copies one after another, each by its own optimal schedule, attains it, because
consecutive copies are tool-disjoint and so the magazine must be rebuilt in full
between them. For $H$ the difficulty is that a group may in principle contain jobs
from two different copies, since two jobs requiring two tools each can together
require exactly $b=4$. Verifying that no such grouping helps is what the computation
above does, and what a proof would have to establish in general.
\end{support}

Finally, the phenomenon behind all of this deserves to be stated as a question, since
it is the natural one to ask next and it has no answer here.

\begin{problem}\label{prob:ktns-closes}
Characterise the instances on which the KTNS step closes the walk gap, that is, on
which $H<\Hwalk$. Equivalently: when does re-optimising the loading over the whole
flattened order recover an optimum that no single magazine per group can reach?
\end{problem}

The evidence assembled here is that this happens very often, and that the walk model
therefore substantially overstates the weakness of the standard method.

\subsection{What the gap does not depend on}\label{ssec:clutter}

The results above are quantitative. This subsection asks from which combinatorial
structure of an instance the gap could, even in principle, be read. The natural
object is the family of feasible groups.

\begin{definition}[Group complex and circuit clutter]\label{def:complex}
Let $\mathcal{F}=\rset{S\subseteq J:\lvert\bigcup_{j\in S}T_j\rvert\le b}$. Since
$\lvert T_j\rvert\le b$ and membership is preserved under taking subsets,
$\mathcal{F}$ is an independence system on $J$. Its inclusion-minimal non-members,
the \emph{circuits}, form a clutter $\mathcal{C}$.
\end{definition}

Three plain ideas sit in that definition. That $\mathcal{F}$ is an independence system
means only that feasibility is preserved downwards: remove a job from a group that
fits, and it still fits. The circuits are the minimal ways to fail: a set of jobs
whose tools overflow the magazine although every proper subset fits. A clutter is a
family no member of which contains another, which the circuits form automatically.

\begin{proposition}[$\Kst$ is a hypergraph chromatic number]\label{prop:chromK}
A set of jobs is a feasible group if and only if it contains no circuit. Hence $\Kst$
is the chromatic number of the hypergraph $(J,\mathcal{C})$. The \emph{conflict
graph}, whose edges join pairs of jobs that cannot share a configuration, consists
exactly of the two-element circuits, so its chromatic number is at most $\Kst$ and
may be strictly smaller. For $T_1=\rset{1,2}$, $T_2=\rset{1,3}$, $T_3=\rset{1,4}$
with $b=3$ there is no conflicting pair at all, yet the three jobs together require
four tools, so they form a circuit and $\Kst=2$.
\end{proposition}

\begin{proof}
An infeasible set contains a minimal infeasible subset and feasibility is preserved
downwards, which gives the first claim. Colour classes of a proper colouring of
$(J,\mathcal{C})$ are then exactly feasible groups. Two-element circuits are precisely
the conflicting pairs, and every proper colouring of the hypergraph properly colours
the graph. The stated instance computes as claimed.
\end{proof}

Figure~\ref{fig:hypergraph} draws the example. The point it makes is that pairwise
conflicts undercount: three jobs can be pairwise compatible and jointly infeasible,
so the true obstruction is the hypergraph of all minimal clashes.

\begin{figure}[ht]\centering
\begin{tikzpicture}[font=\small,>=Stealth,
   jb/.style={draw=cPurple!75,fill=white,circle,semithick,minimum size=1.25cm,
              inner sep=0pt,align=center,font=\scriptsize\bfseries}]
  \node[font=\bfseries] at (1.4,2.55) {Conflict graph (pairs only)};
  \node[jb] (a1) at (1.4,1.4)  {$j_1$\\[-1.5pt]{\tiny$\{1,2\}$}};
  \node[jb] (b1) at (0.4,-0.4) {$j_2$\\[-1.5pt]{\tiny$\{1,3\}$}};
  \node[jb] (c1) at (2.4,-0.4) {$j_3$\\[-1.5pt]{\tiny$\{1,4\}$}};
  \node[align=center,font=\scriptsize,text=black!65] at (1.4,-1.85)
     {every pair shares tool $1$, so any two\\ need $3\le b$ tools: \textbf{no edge}};
  \node[font=\Large\bfseries,text=black!45] at (4.6,0.35) {vs.};
  \node[font=\bfseries] at (7.8,2.55) {Circuit hypergraph};
  \begin{scope}[on background layer]
    \fill[cOrange!16] (7.8,0.2) circle (1.55);
    \draw[cOrange!85,semithick,dashed] (7.8,0.2) circle (1.55);
  \end{scope}
  \node[jb] (a2) at (7.8,1.15)  {$j_1$};
  \node[jb] (b2) at (6.9,-0.45) {$j_2$};
  \node[jb] (c2) at (8.7,-0.45) {$j_3$};
  \node[text=black!60,font=\scriptsize\itshape] at (9.55,1.5) {one circuit};
  \node[align=center,font=\scriptsize,text=black!65] at (7.8,-1.85)
     {all three \emph{together} need\\ $\{1,2,3,4\}=4>b$: \textbf{one hyperedge}, $\Kst=2$};
\end{tikzpicture}
\caption{Why $\Kst$ is read off the hypergraph of all minimal clashes and never off
the graph of pairwise ones. On this instance no two jobs conflict, because any two of
them need only three tools, which fit. All three together need four and overflow the
magazine, so they form a circuit, and that single circuit already forces two groups.}
\label{fig:hypergraph}
\end{figure}

\begin{proposition}[No matroid structure]\label{prop:nonmatroid}
$\mathcal{F}$ is not a matroid in general. For $T_1=\rset{1}$, $T_2=\rset{2}$,
$T_3=\rset{3,4}$ with $b=2$, both $\rset{3}$ and $\rset{1,2}$ are independent, but
neither $\rset{1,3}$ nor $\rset{2,3}$ is, so the exchange axiom fails. Greedy and
matroid-union arguments are therefore unavailable without further hypotheses.
\end{proposition}

\begin{proof}
The relevant unions have sizes $2$, $2$, $3$ and $3$.
\end{proof}

\begin{proposition}[The circuit clutter does not determine the gap]\label{prop:noclutter}
Two instances with $b=4$, $\lvert\U\rvert=7$, four jobs, isomorphic circuit clutters
and equal $\Kst=3$ and $\Zst=3$ can have different gaps:
\[
\begin{aligned}
  I_0:\ &T=\bigl(\rset{1,3,7},\ \rset{2,3,4},\ \rset{3,4,5,6},\ \rset{4,5,6}\bigr),\\
  I_1:\ &T=\bigl(\rset{1,2,4,7},\ \rset{2,5,6,7},\ \rset{3,5},\ \rset{4,6}\bigr).
\end{aligned}
\]
In both, every pair of jobs conflicts except the last two, and no larger set is a
minimal clash. Yet $I_0$ has gap zero and $I_1$ has gap one, in both the walk value
and the heuristic value.
\end{proposition}

\begin{proof}
The pairwise union sizes are checked directly. Exhaustive enumeration over sequences
and over minimum-cardinality groupings gives $\Zst=3$ and $H=\Hwalk=3$ for $I_0$, and
$\Zst=3$ and $H=\Hwalk=4$ for $I_1$.
\end{proof}

No criterion that reads only the circuit clutter, and a fortiori none that reads only
the conflict graph, can therefore decide when the gap vanishes. Tool-level structure
is essential, and that is exactly the level at which Proposition~\ref{prop:hk3}
operates: its overlap sizes $x_{uv}$ are precisely the data the clutter forgets.

On the covering side, however, the clutter is the right object, and blocking duality
makes this precise.

\begin{proposition}[Circuit--group blocking identity]\label{prop:circblocker}
On ground set $J$, the blocker of the circuit clutter, meaning the clutter of its
minimal transversals, is the family of minimal complements of maximal feasible
groups.
\end{proposition}

\begin{proof}
A set meets every circuit if and only if its complement contains no circuit, if and
only if its complement is a feasible group. Minimal transversals therefore correspond
to maximal feasible complements.
\end{proof}

This has a concrete consequence for the next section. Grouping branch-and-price rests
on a set-covering relaxation, and it is fast only when that relaxation is
\emph{ideal}, meaning its fractional optimum is already integral. The next result
exhibits the smallest family where it is not.

\begin{proposition}[Odd rings are not ideal]\label{prop:oddring}
Consider the clutter of job clusters over the maximal feasible groups; its blocker is
the clutter of minimal group-covers, so $\Kst$ is the blocker's minimum size. On the
$k$-ring with $b=3$ and $k\ge4$, the fractional cover equals $k/2$ while
$\Kst=\lceil k/2\rceil$. The covering relaxation is therefore tight for every even
$k$ and has integrality gap $\tfrac12$ for every odd $k$.
\end{proposition}

\begin{proof}
For $k\ge4$ the maximal feasible groups are exactly the pairs of adjacent jobs.
Weight $\tfrac12$ on each covers every job exactly once, so the fractional cover is at
most $k/2$. The dual packing $y_j=\tfrac12$ is feasible, since no group holds more than
two jobs, so by linear-programming duality the fractional cover equals $k/2$. The
integer optimum is $\lceil k/2\rceil$.
\end{proof}

The hypothesis $k\ge4$ is needed: at $k=3$ all three jobs require only the three tools
$\rset{1,2,3}$, so the whole job set is one feasible group and $\Kst=1$.

This locates a first provable obstruction for the set-covering relaxation behind
grouping branch-and-price \citep{Crama1994Column,Otiai2024}: odd circular structure,
in exact parallel to odd holes in perfect-graph theory \citep{Chudnovsky2006}. Whether
odd-ring-like structure is the only obstruction is open, and
Section~\ref{ssec:further} records what is known about it.

\subsection{The setup-cost spectrum}\label{ssec:collapse}

In practice a machine stop costs time beyond the individual tool exchanges. This
subsection shows that the SSP and the JGP are the two ends of a single parameterised
problem, with the transition governed by the gap itself. \citet{Burger2015} pose
exactly this question at the end of their paper.

Charge, in addition to the tool switches, a cost $\rho\ge0$ for every
\emph{configuration change}. The objective becomes
\[
  \text{switches}\;+\;\rho\cdot(\text{number of configuration changes}),
\]
where $\rho$ is the ratio of the fixed stop time to the time of one tool exchange. A
grouping into $p$ configurations makes $p-1$ changes, so the setup term alone is
minimised by the Job Grouping Problem. As $\rho$ grows it dominates and pulls the
optimum onto minimum-cardinality groupings.

\begin{theorem}[Setup-cost collapse]\label{thm:collapse}
If $\rho>\Hwalk-\Zst$, then every optimal solution of the setup-augmented objective
uses a minimum-cardinality grouping, and the two-phase heuristic is exact for that
objective. The collapse threshold therefore satisfies $\rho_c\le\Hwalk-\Zst$, and for
the ring family $\rho_c=1$.
\end{theorem}

\begin{proof}
Every feasible grouping $\mathcal{P}$ has $\lvert\mathcal{P}\rvert\ge\Kst$ and, by
Theorem~\ref{thm:grouping}, $\gamma(\mathcal{P})\ge\Zst$, so its augmented cost is at
least $\rho(\lvert\mathcal{P}\rvert-1)+\Zst$. The best minimum-cardinality grouping
has augmented cost $\rho(\Kst-1)+\Hwalk$. If $\lvert\mathcal{P}\rvert>\Kst$, the
former exceeds the latter whenever
$\rho(\lvert\mathcal{P}\rvert-\Kst)>\Hwalk-\Zst$, which holds for every such
$\mathcal{P}$ as soon as $\rho>\Hwalk-\Zst$. No grouping of larger cardinality is
therefore optimal, and the heuristic, which selects the cheapest ordering of a
minimum-cardinality grouping, attains the optimum.
\end{proof}

The threshold must be read in the walk value. Using the heuristic value instead makes
the statement false, because on the $6$-ring $H-\Zst=0$ while a four-configuration
grouping still wins for small positive $\rho$.

\begin{lemma}[The other end of the spectrum]\label{lem:lowrho}
If $\rho<1/(n-1)$, every optimum of the setup-augmented objective is a
switch-optimal SSP schedule.
\end{lemma}

\begin{proof}
A schedule makes at most $n-1$ configuration changes. Let $S^{*}$ be switch-optimal.
Any schedule with at least $\Zst+1$ switches has augmented cost at least
$\Zst+1>\Zst+\rho(n-1)$, which is at least the augmented cost of $S^{*}$. No such
schedule is therefore optimal.
\end{proof}

Together these pin both ends. For $\rho<1/(n-1)$ the augmented problem is the SSP.
For $\rho>\Hwalk-\Zst$ its optima use minimum-cardinality groupings and the two-phase
heuristic is exact, so the problem has collapsed onto the JGP. The transition is
confined to the interval $[\,1/(n-1),\ \Hwalk-\Zst\,]$, which is to say it is governed
by the gap itself.

\begin{remark}[On practical values of $\rho$]\label{rem:rho}
It is natural to ask where real manufacturing sits on this spectrum. Answering that
requires measured values of the stop time and the per-tool exchange time in a
comparable setting, and no source giving both could be found during this internship.
The spectrum result is therefore stated as it is proved, as a statement about the
parameter $\rho$, and no claim is made about where practice falls on it. Obtaining
such measurements would settle a question that Theorem~\ref{thm:collapse} otherwise
leaves open.
\end{remark}

\subsection{Which grouping to choose}\label{ssec:grouping-sel}

Theorem~\ref{thm:grouping} turns the heuristic's weakness into a concrete algorithmic
question. The optimum is attained by \emph{some} feasible grouping, and
minimum-cardinality ones may miss it (Example~\ref{ex:ring-why}). On which grouping,
then, should the sequencing phase be run?

Admitting a single extra group already recovers the ring optimum in the walk model. In
general the trade-off is between more configuration changes and cheaper individual
transitions. At $\Kst=3$, Proposition~\ref{prop:hk3} makes the criterion explicit:
choose the grouping whose configurations have the smallest pairwise overlap outside
the third. For general $\Kst$, characterising the groupings that attain
$\min_{\mathcal{P}}\gamma(\mathcal{P})$, or giving an efficient rule that improves on
the minimum-cardinality choice, is open.

Two things follow from Section~\ref{ssec:ktnsgap} and are worth noting here. The
first is that the practical value of such a rule is smaller than the walk model
suggests, since the KTNS step already recovers the optimum on most instances where
the walk model loses it. The second is that the instances where a rule would help are
identifiable in advance by the same criterion that identifies positive gaps: a
magazine large relative to the jobs.
